
	\newpage
% =========================================================
	\section{Fixpunkte}
% =========================================================
	
	\defi{
		\n 
		Was ist nochmal ein Fixpunkt? \n 
		Ein Fixpunkt ist ein Punkt $x$ einer Funktion $f$, sodass $f(x) = x$ gilt. Der Punkt bildet über die Funktion auf sich selbst ab.
		\n 
		Einen solchen "Fixpunkt" zu finden ist bei einer beliebigen Funktion $f$ nicht trivial.
	}

	\sbreak
	\subsection{Newton Verfahren}
	
	\defi{
		\n 
		Das Verfahren von Newton kann die Nullstellen von $f(x)$ approximieren und bewerkstelligt dies mithilfe von Tangenten.
		Das Problem: 
		\begin{align*} 
 			f(x) = a \aae x \approx a 
 		\end{align*}
%		Lösung: 
% 			f(x) - a = 0 
 		Herleitung: Approximation von Nullstellen mithilfe von Tangenten $t$ in $x_n$
 		\begin{align*} 
	 		t(x) &= f(x_n) + f'(x_n)\cdot (x-x_n)  \\
	 		t(x_{n+1}) &= 0 \\
	 		f(x_n) + f'(x_n) \cdot (x_{n+1} - x_n) &= 0 \\
	 		x_{n+1} &= x_n - \frac{f(x_n)}{f'(x_n)} 
 		\end{align*}
	}
	
	\methode{
		\n 
		Die Formel für das Newton Verfahren:
		\begin{align} 
 			x_{n+1} = x_n - \frac{f(x_n)}{f'(x_n)} 
 		\end{align}
		Informell bzw. graphisch bedeutet das: Man legt eine Tangente auf den Startwert $x_0$ und bestimmt davon die Nullstelle. Dann geht man an den Punkt der Funktion mit derselben $x$-Koordinate wie die Nullstelle und beginnt von vorne (Tangente anlegen, davon NST ausrechnen etc\dots). \\
%		Ein Video, welches die graphische Lösungsbestimmung veranschaulicht: 
		\n 
		Was hat das mit Fixpunkten zu tun? Naja, wenn man eine Nullstelle in das Newton Verfahren eingibt, kommt derselbe Punkt der Funktion wieder heraus. Nullstellen sind also Fixpunkte der Iterationsvorschrift vom Newton Verfahren.
	}
	
	
	\sbreak
	\bsp{
		\begin{align*} 
 			f(x) = x^3 - 2x^2 - 1 \aae f'(x) = 3x^2 - 4x 
 		\end{align*}
		Der Startpunkt sei $x_0 = 1$. Ich führe drei Schritte von Newton aus:
		\begin{align*} 
 			x_1 &= x_0  - \frac{f(x_0)}{f'(x_0)} \\
 			&= 1 - \frac{1 - 2 - 1}{3 - 4} = -1  \\ \\
 			x_2 &= x_1 - \frac{f(x_1)}{f'(x_1)} \\
			&= (-1) - \frac{(-1) - 2 - 1}{3 + 4} = -\frac{3}{7} \\ \\
			x_3 &= (-\frac{3}{7}) - \frac{(\frac{27}{343}) + (\frac{18}{49}) - 1}{(\frac{27}{49}) + (\frac{12}{7})} = \dots 
 		\end{align*}
		Dieses Verfahren würde sich der Nullstelle bei $x = 2.2056$ annähern. 
		\n
		Keine Sorge, in der Klausur werden nur solche Funktionen und Zahlen gewählt werden, dass ihr alles im Kopf ausrechnen könnt (also hier ist $x_3$ deutlich zu schwer im Kopf auszurechnen).
	}
	
	\newpage
% =========================================================	
	\subsection{Eigenschaften}
% =========================================================

	\defi{
		\n 
		Allgemein versuchen Fixpunktverfahren (wie eben Newton) mit einem Fixpunkt $x*$, der Iterationsvorschrift $\Phi(x)$ und einem Startwert $x_0$
		\begin{align} 
 			x* = \Phi(x*) 
 		\end{align}
		mit Gleichungen der folgenden Form zu lösen:
		\begin{align} 
			x_{i+1} = \Phi(x_i) 
 		\end{align}
		Aber wo sind grafisch von der Iterationsfunktion $\Phi(x)$ die Fixpunkte der Funktion $f(x)$? \\
		Man könnte es sich nicht einfacher vorstellen: Die Fixpunkte sind die Schnittpunkte der Graphen von $\Phi(x)$ und $f(x)$.
		\n
		Die Iterationen kann man graphisch auch schön veranschaulichen: man bildet eine Treppe und springt zwischen Iterationsfunktion und der eigentlichen Funktion $f$ hin und her. Also man fängt auf der Funktion am Startwert an und zieht eine gerade Linie nach links bzw. rechts, bis man auf $\Phi(x)$ trifft. Dann zieht man eine gerade Linie nach oben bzw. unten, bis man wieder auf $f(x)$ landet. Das wäre dann eine Iteration.
		\n 
		Also gerade Linien nach oben bzw. unten zeichnen, wenn man von $\Phi(x)$ kommt und wieder nach rechts bzw. links "abbiegen", wenn man von $f(x)$ kommt.
		\sbreak \n 
		Es gibt außerdem verschiedene Arten von Fixpunkten (von $\Phi$ abhängig!):
		\begin{itemize}
			\item $\abs{\Phi'(x*)} < 1 \Longrightarrow$ anziehender Fixpunkt (konvergiert)
			\item $\abs{\Phi'(x*)} > 1 \Longrightarrow$ abstoßender Fixpunkt (divergiert)
			\item $\abs{\Phi'(x*)} = 1 \Longrightarrow$ indifferent
		\end{itemize}
		Anziehend bedeutet im Endeffekt (informell), dass Verfahren eine Tendenz haben, auf diesen Fixpunkt zuzugehen. Abstoßend ist hierbei verständlicherweise das Gegenteil. Grafisch würde die Treppe bei dem Fixpunkt konvergieren (anziehender Fixpunkt) oder davon weg divergieren (abstoßender Fixpunkt). 
		\n
		%Hier wird es schön veranschaulicht: \\
	}

	\newpage
	\bsp{
		\n 
		Es sei die Iterationsvorschrift $\Phi(x) = x^2$ gegeben. Wir sollen alle Fixpunkte bestimmen und zeigen, ob diese anziehend oder abstoßend sind.
		\n 
		Fixpunkte bestimmen:
		\begin{align*} 
 			x &= x^2 \\
 			0 &= x^2 - x \\
 			\text{p-q-Formel: }\; x_{0,1} &= - \frac{-1}{2} \pm \sqrt{\frac{1}{4} - 0} \\
 				&= \frac{1}{2} \pm \frac{1}{2}
 		\end{align*}
		Die Fixpunkte sind $x_0 = 0$ und $x_1 = 1$.
		\n 
		Welcher Art gehören die Fixpunkte an?
		\begin{align*} 
 			\Phi'(x) &= 2x \\
 			\Phi'(0) &= 0 < 1 \\
 			\Phi'(1) &= 2 > 1 
 		\end{align*}
		Ergo ist $x_0 = 0$ ein anziehender Fixpunkt und $x_1 = 1$ ein abstoßender!
	}

	\sbreak
	\defi{
		\n
		Im Normfall konvergiert das Newton-Verfahren mit quadratischer Konvergenz. Bei mehrfachen Nullstellen besteht aber nur noch lineare Konvergenz. Wir können dieses Problem vorbeugen, indem wir das Näherungsverfahren modifizieren.
		\n 
		Bei einer $k$-fachen Nullstelle modifiziert man das newton'sche Näherungsverfahren mit einem Faktor $k$:
		\begin{align}
			x_{n+1} = x_n - k \cdot \frac{f(x_n)}{f'(x_n)} 
		\end{align}
	}


	\newpage
	\defi{ Wdh. mehrfache Nullstellen:
		\n 
		Falls jemand nicht weiß, was eine mehrfache Nullstelle $x_0$ von $f$ ist, hier eine kurze Erklärung:
		\n
		Für jedes mal, dass man $f$ ableiten kann und $x_0$ immer noch eine Nullstelle ist, ist die Vielfachheit dieser Nullstelle $x_0$ um eins erhöht.
		\n 
		Formal: \\
		Eine Nullstelle $x_0$ ist eine $k$-fache Nullstelle, wenn
		\begin{align}
			f^{(j)}(x_0) = 0 \aae \forall j\in [0, k-1]
		\end{align}
		Sei $f$ nun mindestens $k$-mal differenzierbar. Ist $x_0$ eine $k$-fache Nullstelle, aber keine $(k + 1)$-fache, so ist $k$ die Vielfachheit der Nullstelle $x_0$.
	}

	\sbreak
	\bsp{
		\n 
		Gegeben sei
		\begin{align*}
			f(x) &= x^2 - 2x + 1 \\
			x_0 &= 1
		\end{align*}
		$x_0$ ist eine doppelte Nullstelle, da
		\begin{align*}
			f(1) &= 1^2 - 2 + 1 = 0 \\
			f'(x) &= 2x - 2 \\
			f'(1) &= 2 - 2 = 0
		\end{align*}
		$x_0$ ist keine dreifache Nullstelle, weil
		\begin{align*}
			f"(x) &= 2 \\
			f"(1) &= 2 \ne 0
		\end{align*}
	}

	\sbreak
	\wichtig{
		\n
		Wir sagen, dass ein Iterationsverfahren $\Phi(x)$ genau \fat{dann gut} ist, wenn die Lösung des Iterationsverfahren ein \fat{anziehender Fixpunkt} ist. 
		\\ Ansonsten wird der Fixpunkt halt nie durch das Verfahren erreicht oder es ist nicht einmal ein Fixpunkt, das wäre natürlich Quatsch!
	}


	\newpage
	\bsp{
		\n 
		Wir sollen zeigen, ob folgende Iterationsvorschriften sich zum Ausrechnen von Nullstellen einer beliebigen Funktion $f(x)$ eignen:
		\begin{itemize}
			\item Gerade, die nach oben verschoben ist:
				\begin{align*} 
 					\Phi_0(x) = x + 1  
 				\end{align*}
				Wir schauen nach, ob Nullstellen Fixpunkte der Iterationsfunktion sind:
				\begin{align*}
 					\Phi_0(0) = 0 + 1 = 1 
 				\end{align*}
				Das ist nicht der Fall und somit eignet sich diese Vorschrift nicht für das Ausrechnen von Nullstellen!
			\item Eine Gerade durch den Ursprung: 
				\begin{align*} 
 					\Phi_1(x) = 2x  
 				\end{align*}
				\begin{align*} 
 					\Phi_1(0) = 0  
 				\end{align*}
				Null ist ein Fixpunkt der Vorschrift, ist dieser auch anziehend?
				\begin{align*} 
 					\abs{\Phi_1'(0)} = 2 > 1 
 				\end{align*}
				$\Longrightarrow$ abstoßend $\Longrightarrow$ eignet sich nicht
			\item Potenzfunktion:
				\begin{align*} 
 					\Phi_2(x) = x^3 
 				\end{align*}
				\begin{align*} 
 					\Phi_2(0) = 0 
				\end{align*}
				Null ist ein Fixpunkt, jetzt noch Eigenschaft des Fixpunktes überprüfen:
				\begin{align*} 
 					\abs{\Phi_2'(0)} = 3 \cdot 0^2 = 0 < 1 
 				\end{align*}
				$0$ ist ein anziehender Fixpunkt von $\Phi_2(x)$ und somit eignet sich die Vorschrift hierfür (da die Iterationsfunktion aber gar nicht von $f(x)$ abhängig ist lässt sich darüber streiten, wie sinnvoll das wirklich ist).
		\end{itemize}
	}


	\newpage
% =========================================================	
	\subsection{Banach'scher Fixpunktsatz}

	\defi{
		\n 
		Der Banach'sche Fixpunktsatz sagt folgendes aus:
		\n 
		Es existiert genau ein anziehender Fixpunkt in einem Intervall $I$ genau dann, wenn
		\begin{enumerate}
			\item $I$ abgeschlossen
			\item $\forall x \in I: \; \Phi(x) \in I$
			\item $\Phi$ ist "Lipschitz-stetig" im Intervall $I$ mit $0 < L < 1$:
				\begin{align} 
 					\exists L: \abs{\Phi(x) - \Phi(y)} \le L \cdot \abs{x - y} 
 				\end{align}
		\end{enumerate}
		Wenn die Bedingungen erfüllt sind, hat das auch noch eine andere Folgerung. Nicht nur ist in dem Intervall $I$ dann genau ein anziehender Fixpunkt, sondern egal welchen Startwert $x_0 \in I$ man wählt, die Fixpunktiteration $\Phi(x)$ würde auf ebenjenen anziehenden Fixpunkt garantiert konvergieren.
	}
	
	% Vereinfachung:
	% Gelte	∀x∈I: |Φ′(x)| < 1  (L:= max x∈I|Φ′(x)|)
	% When true:
	% Dann konvergiert das Iterationsverfahren für ein 	x0∈I gegen den eindeutigen Fixpunkt x∈I
	
	\sbreak
	\methode{
		\begin{enumerate}
			\item Ein Intervall $I$ ist abgeschlossen, wenn es beide Grenzen enthält. \\
			$0 \le x \le 1$ wäre abgeschlossen, \\
			$0 < x < 1$ ist ein offenes Intervall.
			\item Herausfinden, wann in dem Intervall die Funktionswerte am kleinsten und am größten sind, einsetzen und schauen, ob sich diese noch im Intervall befinden (Weil alle Werte dazwischen müssen sich dann auch im Intervall befinden) \\
			Gute Startüberlegung wäre "monoton steigend/fallend" und ähnliches.
			\item Für diesen Punkt verwenden wir immer den Mittelwertsatz (MWS):
			\begin{align} 
 				\exists \xi \in [a, b]: f(b) - f(a) = f'(\xi) \cdot (b - a) 
 			\end{align}
			Wenn man dies jetzt auf $\Phi$ anwendet gilt dann:
			\begin{align} 
 				\abs{\Phi(x) - \Phi(y)} \le \abs{x -y} \cdot \max_{z \in [x, y]} \abs{\Phi'(z)} 
 			\end{align}
			Das gilt daher, weil man durch den Mittelwertsatz weiß, dass die zwei Gleichungen für ein bestimmtes $\xi$ gleich ist und für den maximalen Wert der Ableitung im Intervall $I$ also größer oder gleich der linken Seite sein muss.
		\end{enumerate}
	}
		
	\newpage
	\bsp{
		\n 
		Es sei die Iterationsvorschrift $\Phi(x) = x^2$ vom obigem Beispiel gegeben. \\
		Jetzt sollen wir mithilfe vom Fixpunktsatz beweisen, dass 0 ein anziehender Fixpunkt ist (bzw. wir zeigen, dass in einem Intervall, wo auch die 0 drin ist, ein anziehender Fixpunkt existiert!).
		\n
		Sei 
		\begin{align*}
			I = \l[-\frac{1}{4};\; \frac{1}{4} \r]
		\end{align*}
		\begin{enumerate}
			\item $I$ ist abgeschlossen, da trivial (bzw. das Intervall hat keine runde Klammern mit Inhalten ungleich Unendlich)
			\item $\Phi(x)$ ist im Bereich von $[-0.25; 0]$ und $[0; 0.25]$ monoton fallend/steigend (Ableitung ist $2x$), daraus wissen wir, dass der kleinste/größte Funktionswert bei $x = 0$ bzw. $x = \pm 0.25$ liegen muss.
			\begin{align*}
				\Phi(-0.25) &= 0.125 \\
				\Phi(0) &= 0
			\end{align*}
			Diese Funktionswerte liegen immer noch im Intervall $I$.
			\item Auf $\Phi$ angewendet gilt:
			\begin{align*} 
 				\abs{\Phi(x) - \Phi(y)} \le \abs{x -y} \cdot \max_{z \in [x, y]} \abs{\Phi'(z)} 
 			\end{align*}
			Wir wissen, dass $\Phi'(z) = 2z$ ist und der maximale Wert dieser Funktion in unserem Intervall $I$ sich an der Stelle $z = 0.25$ befindet (die Ableitung ist streng monoton steigend).
			\begin{align*} 
 				\abs{\Phi(x) - \Phi(y)} \le \max_{z \in I}(2z) \cdot \abs{x - y} = \frac{1}{2} \cdot \abs{x -y} 
 			\end{align*}
			Die "Lipschitz-stetigkeit" ist also mit $L = \frac{1}{2}$ erfüllt.
			\n 
			Die Fixpunktiteration konvergiert also für alle Startwerte aus dem Intervall $I:= [-0.25; 0.25]$ gegen den Fixpunkt $x_0 = 0$!
		\end{enumerate}
	}
